\bigskip\noindent{\Large\bfseries\color{apren} Test 3. Test global de nombres complexos}\par\medskip
\iniciTest{complexosT3}{c,c,b,c,d,a,a,c,d,b}
\begin{enumerate}
\pregunta Considerem els complexos següents:
\[
\begin{array}{l}
z_{1}=2-3i \\
z_{2}=-5+i \\
z_{3}=-1-i
\end{array}
\]
aleshores:
\begin{enumerate}
\opcio{a}{\[
z_{1}\cdot \overline{z_{3}}+5z_{2}=24+10i
\]}
\opcio{b}{\[
\frac{z_{2}}{z_{3}}+\overline{z_{1}}=4i
\]}
\opcio{c}{\[
z_{1}+z_{2}^{2}+z_{3}^{3}=28-15i
\]}
\opcio{d}{Cap de les anteriors}
\end{enumerate}\feedback
\pregunta El quocient
\[
\frac{3-2ai}{4+3i}
\]
és:
\begin{enumerate}
\opcio{a}{un nombre real si $a=2$}
\opcio{b}{un imaginari pur si $a=\frac{9}{8}$}
\opcio{c}{és $3+3i$ si $a=-\frac{21}{2}$}
\opcio{d}{Cap de les anteriors}
\end{enumerate}\feedback
\pregunta Quina de les afirmacions següents és falsa?
\begin{enumerate}
\opcio{a}{$x=\frac{12}{5}-\frac{31}{5}i,y=\frac{7}{5}+\frac{9}{5}i$ és solució del sistema
\[
\left\{
\begin{array}{l}
2x+2(1+i)y=4-6i \\
x+3iy=-3-2i
\end{array}
\right.
\]}
\opcio{b}{$x^{4}-4x^{3}+6x^{2}-4x+5$ és un polinomi que té les arrels $i
$ i $1-i$}
\opcio{c}{\[
\frac{i^{316}-i^{205}}{i^{30}+i^{187}}=i
\]}
\opcio{d}{$\sqrt{2}\pm \sqrt{3}i$ són les solucions complexes de l’equació
\[
x^{2}-2\sqrt{2}x+5=0
\]}
\end{enumerate}\feedback
\pregunta Quina de les afirmacions següents és correcta?
\begin{enumerate}
\opcio{a}{L’expressió en forma polar del complex conjugat de $2+2\sqrt{3}i
$ és $4_{\frac{\pi }{6}}$}
\opcio{b}{L’expressió en forma trigonomètrica de l’invers del complex $
-3+\sqrt{3}i$ és
\[
\frac{\sqrt{3}}{6}\left( \cos 150
^{\circ}
+i\sin 150
^{\circ}
\right)
\]}
\opcio{c}{Si $z=4_{\frac{\pi }{4}}$, aleshores l’expressió en forma binòmica del complex $-\overline{z}$ és
\[
-2\sqrt{2}+2\sqrt{2}i
\]}
\opcio{d}{Cap de les anteriors}
\end{enumerate}\feedback
\pregunta Quina de les afirmacions següents és falsa?
\begin{enumerate}
\opcio{a}{Perquè
\[
\frac{a+2i}{3+bi}=\left( \sqrt{2}\right) _{\frac{\pi }{4}}
\]
és $a=4$ i $b=-1$}
\opcio{b}{L’equació
\[
x^{2}-3x+9=0
\]
té com solucions $3_{\frac{\pi }{3}}$ i $3(\cos 300
^{\circ}
+i\sin 300
^{\circ}
)$}
\opcio{c}{Si
\[
z=\left( \cos \frac{\pi }{13}+i\sin \frac{12\pi }{13}\right) ^{13}
\]
aleshores $\left\vert z\right\vert =1$ i $\arg z=\pi $}
\opcio{d}{Si $z$ és un complex no nul, aleshores es compleix
\[
\left\vert \frac{z^{-1}}{\overline{z}}\right\vert =\frac{1}{\left\vert
z\right\vert ^{2}}\text{ \ \ i \ \ }\arg \left( \frac{z^{-1}}{\overline{z}}
\right) =2\arg z
\]}
\end{enumerate}\feedback
\pregunta Considerem els complexos següents
\[
\begin{array}{l}
z_{1}=2(\cos 15
^{\circ}
+i\sin 15
^{\circ}
) \\
z_{2}=3_{\frac{\pi }{4}} \\
z_{3}=-1-i
\end{array}
\]
aleshores, quin dels següents resultats és fals?
\begin{enumerate}
\opcio{a}{\[
z_{1}^{2}\cdot z_{2}^{4}=162\sqrt{3}-162i
\]}
\opcio{b}{\[
z_{1}^{4}+z_{2}^{2}+z_{3}^{3}=10+(7+8\sqrt{3})i
\]}
\opcio{c}{\[
\left( \frac{z_{1}}{z_{3}}\right) ^{4}\cdot z_{2}^{2}=18\sqrt{3}-18i
\]}
\opcio{d}{\[
\left( \frac{z_{2}\cdot z_{3}}{z_{1}}\right) ^{6}=\frac{729}{8}i
\]}
\end{enumerate}\feedback
\pregunta Se sap que l’equació
\[
z^{3}+az^{2}+6(1-i)z+b=0
\]
amb $a,b\in \mathbb{R}$ té la solució $1+i$, aleshores
\begin{enumerate}
\opcio{a}{$a=-1,b=-10$ i les altres dues solucions de l’equació són $1+2i$
i $-1-3i$}
\opcio{b}{$a=1,b=-10$ i les altres dues solucions de l’equació són $1-i$ i $
-1$}
\opcio{c}{$a=1,b=10$ i les altres dues solucions de l’equació són $1-i$ i $1
$}
\opcio{d}{$a=-1,b=10$ i les altres dues solucions de l’equació són $1-i$ i $
2$}
\end{enumerate}\feedback
\pregunta Quina de les igualtats següents és certa?
\begin{enumerate}
\opcio{a}{\[
\sqrt{\frac{1}{2}+\frac{\sqrt{3}}{2}i}=\left\{
\begin{array}{l}
\frac{1}{2}+\frac{\sqrt{3}}{2}i \\
-\frac{1}{2}-\frac{\sqrt{3}}{2}i
\end{array}
\right.
\]}
\opcio{b}{\[
\sqrt[3]{\sqrt{3}+i}=\left\{
\begin{array}{l}
\left( \sqrt[3]{2}\right) _{30
^{\circ}
} \\
\left( \sqrt[3]{2}\right) _{150
^{\circ}
} \\
\left( \sqrt[3]{2}\right) _{270
^{\circ}
}
\end{array}
\right.
\]}
\opcio{c}{\[
\sqrt[3]{\frac{3_{30
^{\circ}
}}{\left( \frac{1}{9}\right) _{60
^{\circ}
}}}=\left\{
\begin{array}{l}
3_{110
^{\circ}
} \\
3_{230
^{\circ}
} \\
3_{350
^{\circ}
}
\end{array}
\right.
\]}
\opcio{d}{Cap de les anteriors}
\end{enumerate}\feedback
\pregunta Quin dels complexos següents satisfà la
condició que el cub del seu conjugat coincideix amb el seu oposat?
\begin{enumerate}
\opcio{a}{$\cos 30
^{\circ}
-i\sin 30
^{\circ}
$}
\opcio{b}{$\cos 90
^{\circ}
-i\sin 90
^{\circ}
$}
\opcio{c}{$\cos 60
^{\circ}
+i\sin 60
^{\circ}
$}
\opcio{d}{$\cos 135
^{\circ}
+i\sin 135
^{\circ}
$}
\end{enumerate}\feedback
\pregunta Quina de les afirmacions següents és falsa?
\begin{enumerate}
\opcio{a}{Si un triangle té els seus vèrtexs en els punts $(2,0),(8,1)$
i $(4,5)$, en girar-lo un angle recte al voltant de l’origen s’obté
un altre triangle de vèrtexs $(0,2),(-1,8)$ i $(-5,4)$}
\opcio{b}{Si $(2,7)$ i $(6,5)$ són dos vèrtexs oposats d’un quadrat, les
coordenades dels altres dos vèrtexs són $(3,4)$ i $(8,5)$}
\opcio{c}{Si $(1,1)$ és un vèrtex d’un triangle equilàter el qual
baricentre és en $(0,3)$, aleshores les coordenades d’un altre dels seus vèrtexs són
\[
\left( -\frac{1}{2}+\sqrt{3},4+\frac{\sqrt{3}}{2}\right)
\]}
\opcio{d}{Si $(0,1)$ és un vèrtex d’un hexàgon regular centrat en el
origen, aleshores les coordenades d’un altre dels seus vèrtexs són
\[
\left( \frac{\sqrt{3}}{2},\frac{1}{2}\right)
\]}
\end{enumerate}\feedback
\end{enumerate}
